\subsubsection{迹定理}\label{subsec:sobolev-trace-theorem}

本小节讨论 Sobolev 空间中的函数能否限制到区域的边界, 或更一般地限制到其他曲面的问题. 也关注其逆问题, 即边界数据必须多光滑, 才能由 $H^k(\Omega)$ 中的函数取到这些边界数据. 首先考察定义在 $\mathbb{R}^m$ 上的函数及其在 $\mathbb{R}^{m-1}\times\{0\}$ 上的限制.

\begin{Theorem}
\label{thm:hyperplane-sobolev-trace}
设实数 $s>1/2$. 那么存在连续线性映射 $T:H^s(\mathbb{R}^m)\to H^{s-1/2}(\mathbb{R}^{m-1})$, 称为迹算子, 它具有以下性质: 对任意 $\phi\in\mathcal{D}(\mathbb{R}^m)$, 有
\begin{equation}\label{eq:hyperplane-trace-test-function-restriction}
T\phi(x_1,\ldots,x_{m-1})=\phi(x_1,x_2,\ldots,x_{m-1},0).
\end{equation}

\end{Theorem}

可以证明, 对 $s\leq1/2$, 该结果不成立; 事实上, 此时在 $x_m=0$ 的一个邻域内为零的函数在 $H^s(\mathbb{R}^m)$ 中稠密.

\begingroup
\postdisplaypenalty=10000
\renewcommand{\endProof}{\par\nobreak\hfill$\blacksquare$\par\vspace{0.1cm}}
\begin{Proof}
设 $\phi\in\mathcal{D}(\mathbb{R}^m)$, 并令 $g(\bm{x}^{\prime})=\phi(\bm{x}^{\prime},0)$. 用 $\widetilde\phi$ 表示 $\phi$ 仅对第 $m$ 个变量的 Fourier 变换, 即
\begin{equation}\label{eq:trace-proof-partial-normal-fourier-transform}
\widetilde\phi(\bm{x}^{\prime},\xi_m)=\dfrac{1}{\sqrt{2\pi}}\int_{-\infty}^{\infty}\phi(\bm{x}^{\prime},x_m)e^{-ix_m\xi_m}\,\mathrm{d}x_m,
\end{equation}
用 $\hat\phi$ 和 $\hat g$ 分别表示 $\phi$ 和 $g$ 在 $\mathbb{R}^m$ 及 $\mathbb{R}^{m-1}$ 中的 Fourier 变换. Fourier 反演公式给出
\begin{equation}\label{eq:trace-proof-normal-fourier-inversion}
g(\bm{x}^{\prime})=\phi(\bm{x}^{\prime},0)=\dfrac{1}{\sqrt{2\pi}}\int_{-\infty}^{\infty}\widetilde\phi(\bm{x}^{\prime},\xi_m)\,\mathrm{d}\xi_m.
\end{equation}
应用 Fourier 变换, 得到
\begin{equation}\label{eq:trace-proof-tangential-fourier-integral}
\hat g(\bm{\xi}^{\prime})=\dfrac{1}{\sqrt{2\pi}}\int_{-\infty}^{\infty}\hat\phi(\bm{\xi})\,\mathrm{d}\xi_m.
\end{equation}
现在估计
\begin{equation}\label{eq:trace-proof-weighted-fourier-cauchy-estimate}
\begin{aligned}
\|g\|_{s-1/2,2}^2&\leq C_1\int_{\mathbb{R}^{m-1}}|\hat g(\bm{\xi}^{\prime})|^2(1+|\bm{\xi}^{\prime}|^2)^{s-1/2}\,\mathrm{d}\bm{\xi}^{\prime}\\
&=\dfrac{C_1}{2\pi}\int_{\mathbb{R}^{m-1}}\left|\int_{-\infty}^{\infty}\hat\phi(\bm{\xi})\,\mathrm{d}\xi_m\right|^2(1+|\bm{\xi}^{\prime}|^2)^{s-1/2}\,\mathrm{d}\bm{\xi}^{\prime}\\
&\leq C_2\int_{\mathbb{R}^{m-1}}(1+|\bm{\xi}^{\prime}|^2)^{s-1/2}\int_{-\infty}^{\infty}|\hat\phi(\bm{\xi})|^2(1+|\bm{\xi}|^2)^s\,\mathrm{d}\xi_m\\
&\qquad\times\int_{-\infty}^{\infty}(1+|\bm{\xi}|^2)^{-s}\,\mathrm{d}\xi_m\,\mathrm{d}\bm{\xi}^{\prime}.
\end{aligned}
\end{equation}
若 $s>1/2$, 则有
\begin{equation}\label{eq:trace-proof-integrated-normal-fourier-weight}
\begin{aligned}
\int_{-\infty}^{\infty}(1+|\bm{\xi}|^2)^{-s}\,\mathrm{d}\xi_m
&=\int_{-\infty}^{\infty}(1+|\bm{\xi}^{\prime}|^2+\xi_m^2)^{-s}\,\mathrm{d}\xi_m\\
&=(1+|\bm{\xi}^{\prime}|^2)^{-s+1/2}\int_{-\infty}^{\infty}(1+y^2)^{-s}\,\mathrm{d}y.
\end{aligned}
\end{equation}
将其代入 \eqref{eq:trace-proof-weighted-fourier-cauchy-estimate}, 得到对每个 $\phi\in\mathcal{D}(\mathbb{R}^m)$, 有 $\|T\phi\|_{s-1/2,2}\leq C\|\phi\|_{s,2}$. 因而可由连续性将 $T$ 延拓到整个 $H^s(\mathbb{R}^m)$.
\end{Proof}
\endgroup

\begin{Lemma}
\label{lem:continuous-function-trace-and-support}
若 $u\in H^s(\mathbb{R}^m)\cap C(\mathbb{R}^m)$, $s>1/2$, 则 $Tu(\bm{x}^{\prime})=u(\bm{x}^{\prime},0)$. 此外, 对每个 $u\in H^s(\mathbb{R}^m)$, 有 $\operatorname{supp}Tu\subset\operatorname{supp}u\cap(\mathbb{R}^{m-1}\times\{0\})$.
\end{Lemma}

证明留作习题.

现在自然会问, $H^{s-1/2}(\mathbb{R}^{m-1})$ 中哪些元素可作为 $H^s(\mathbb{R}^m)$ 中函数的限制, 即 ``迹'' 得到? 答案是 $H^{s-1/2}(\mathbb{R}^{m-1})$ 中的所有元素都可以用这种方式得到.
